\begin{textsample}{POS Dim 7 – global\_north\_2025\_01 – Score 2.00 – global\_north\_2025\_01/119716144.txt}  \label{ex:f7_pos_396}
Ireland joins South Africa ' s ICJ Gaza genocide case against Israel South Africa started ICJ proceedings against Israel in December 2023 ( Image : PA ) IRELAND has formally joined South Africa ' s genocide case against Israel at the International Court of Justice ( ICJ )."Ireland, invoking Article 63 of the Statute of the Court, filed in the Registry of the Court a declaration of intervention in the case concerning \textbf{Application} of the Convention on the Prevention and Punishment of the Crime of Genocide in the Gaza Strip,"or South Africa versus Israel, the Court said in a statement."There has been a collective punishment of the Palestinian people through the intent and impact of military actions of Israel in Gaza, leaving 44,000 dead and millions of civilians displaced,"The country ' s deputy premier and foreign affairs minister Micheal Martin said at the time. ( Image : PA ) Speaking after a Cabinet meeting, Martin ( above ) added :"We are concerned that a very narrow @ @ @ @ @ @ @ @ @ @ impunity in which the protection of civilians is minimised."Ireland ' s view of the convention is broader and prioritises the protection of civilian life -- as a committed supporter of the convention, the Government will promote that interpretation in its intervention in this case."Intervening in both cases demonstrates the consistency of Ireland ' s approach to the interpretation and \textbf{application} of the Genocide Convention."In December 2023, South Africa started proceedings against Israel, claiming violations of the Genocide Convention in relation to Palestinians in Gaza.
Several countries have since joined the case, including Nicaragua, Colombia, Libya, Mexico, Palestine, Spain and Turkey.
Israel has continued its brutal war on Gaza since a Hamas attack in October 2023 despite a UN Security Council \textbf{resolution} demanding an immediate ceasefire.
Nearly 46,000 people, mostly women and children, have since been killed and over 105,000 injured, according to local health authorities.
Why are you making commenting on The National only available to subscribers?
We @ @ @ @ @ @ @ @ @ @ debate, argue and go back and forth in the comments section of our stories.
We ' ve got the most informed readers in Scotland, asking each other the big questions about the future of our country.
Unfortunately, though, these important debates are being spoiled by a vocal minority of trolls who aren't really interested in the issues, try to derail the conversations, register under fake names, and post vile abuse.
So that ' s why we ' ve decided to make the ability to comment only available to our paying subscribers.
That way, all the trolls who post abuse on our website will have to pay if they want to join the debate -- and risk a permanent ban from the account that they subscribe with.
The conversation will go back to what it should be about -- people who care passionately about the issues, but disagree constructively on what we should do about them.
Let ' s get that debate started!
Callum Baird, Editor of The National Comments : Our @ @ @ @ @ @ @ @ @ @ and valuable part of our community - a place where readers can debate and engage with the most important local issues.
The ability to comment on our stories is a privilege, not a right, however, and that privilege may be withdrawn if it is abused or misused. follow us This website and associated newspapers adhere to the Independent Press Standards Organisation ' s Editors ' Code of Practice.
If you have a complaint about the editorial content which relates to inaccuracy or intrusion, then please contact the editor here.
If you are dissatisfied with the response provided you can contact IPSO here

% matched lemmas: application, resolution
\end{textsample}
